\documentclass[11pt]{article}
\usepackage[T1]{fontenc}
\usepackage{lmodern,amsmath,amssymb,amsthm,booktabs,microtype}
\usepackage[letterpaper,left=1.03in,right=1.03in,top=0.88in,bottom=0.92in]{geometry}
\usepackage{xcolor,hyperref}
\definecolor{ink}{HTML}{244766}
\hypersetup{colorlinks=true,urlcolor=ink,linkcolor=ink,citecolor=ink,
 pdftitle={Recovering omitted primes from normalized zero spacings}}
\usepackage{fancyhdr}
\pagestyle{fancy}\fancyhf{}
\fancyhead[L]{\small\itshape Recovering omitted primes from normalized zero spacings}
\fancyfoot[C]{\small\thepage}
\renewcommand{\headrulewidth}{0.3pt}
\setlength{\headheight}{14pt}
\setlength{\parindent}{1.2em}
\newtheorem{theorem}{Theorem}
\newtheorem{proposition}[theorem]{Proposition}
\newtheorem{corollary}[theorem]{Corollary}
\newcommand{\Z}{\mathbb Z}\newcommand{\Q}{\mathbb Q}
\newcommand{\R}{\mathbb R}
\newcommand{\Prob}{\mathbb P}
\newcommand{\card}{\operatorname{card}}
\renewcommand{\qedsymbol}{$\square$}
\raggedbottom
\begin{document}
\thispagestyle{plain}
\begin{center}
{\large\scshape Analytic number theory}\vspace{0.45em}\\
{\LARGE\bfseries Recovering omitted primes from\\[3pt]normalized zero spacings}\vspace{0.65em}\\
{\normalsize Research note $\cdot$ 26 November 2026}
\end{center}
\begin{abstract}
The zeros of an imprimitive Dirichlet $L$-function on the positive imaginary axis
come from Euler factors omitted when its primitive character is induced to a
larger modulus. Take the consecutive gaps between these zeros, divide by their mean,
and retain only the limiting distribution. If at least three primes were omitted,
the distribution determines their identities. In fact its restriction to
arbitrarily short gaps suffices. The proof combines the known spacing law for
a union of incommensurable lattices with the six exponentials theorem. The inverse
is discontinuous. Disjoint sets of arbitrarily large primes can approximate
the same limiting distribution.
\end{abstract}

\section{What survives normalization}
Let $\chi$ modulo $q$ be induced by the primitive character $\chi^*$ of conductor $f$.
Write
\[
 P=\{p\text{ prime}:p\mid q,\ p\nmid f\},\qquad m=|P|.
\]
Take $P$ nonempty. Euler products in $\Re s>1$ and meromorphic continuation give
\begin{equation}\label{eq:euler}
 L(s,\chi)=L(s,\chi^*)\prod_{p\in P}\bigl(1-\chi^*(p)p^{-s}\bigr).
\end{equation}
At $s=it$ with $t>0$, the primitive functional equation relates $L(it,\chi^*)$
to $L(1-it,\overline{\chi^*})$ by a finite nonzero factor. Classical nonvanishing
on $\Re s=1$ shows that every zero on the positive imaginary axis comes from the finite product
\cite{ked,dlmf}. This includes the primitive character of conductor one, with
$L(s,\chi^*)=\zeta(s)$. Identity~\eqref{eq:euler} is also the change-of-level step
at the end of the Lean proof of Dirichlet GRH \cite{lean}. The boundary-zero
argument uses only classical nonvanishing.

List all positive ordinates of zeros on $\Re s=0$ as $t_1<t_2<\cdots$.
Their intensity (number per unit height) is
\begin{equation}\label{eq:intensity}
 \Lambda=\frac{1}{2\pi}\sum_{p\in P}\log p.
\end{equation}
Thus if absolute height units survive, factoring $\exp(2\pi\Lambda)$ already recovers
$P$. We examine data with this calibration removed. Define the normalized spacings
$g_j=\Lambda(t_{j+1}-t_j)$ and their limiting survival function
\begin{equation}\label{eq:surv}
 S_P(x)=\lim_{N\to\infty}\frac{1}{N}\#\{1\le j\le N:g_j>x\},\qquad x\ge0.
\end{equation}
The $g_j$ have limiting mean one. Common rescaling of all heights leaves them unchanged.

\begin{theorem}[Recovery from short gaps]\label{thm:main}
Let $\chi$ and $\widetilde\chi$ have nonempty omitted-prime sets $P$ and $Q$.
Suppose $S_P(x)=S_Q(x)$ on $0<x<\varepsilon$ for some $\varepsilon>0$.
Then $|P|=|Q|$. If $|P|\ge3$, then $P=Q$. This requires neither an absolute
zero density nor knowledge of the moduli, conductors, or character values.
\end{theorem}

\section{The forward law}
Write $\chi^*(p)=\exp(2\pi i\alpha_p)$ with $0\le\alpha_p<1$. The phase
$\alpha_p$ is rational because a nonzero Dirichlet character value has finite order.
The ordinates contributed by $p$ form the positive part of
\[
  (\Z+\alpha_p)/\lambda_p,\qquad
  \lambda_p=\frac{\log p}{2\pi}.
\]
These lattices are disjoint at positive height. A common $t>0$ for $p$ and $r$
would give $\log p/\log r\in\Q$ from the two nonzero rational numbers
$t\lambda_p$ and $t\lambda_r$, contrary to unique factorization.
Each factor has simple zeros since its derivative equals $\log p$ at a zero.
Counting points below height $T$ now gives $\Lambda T+O(m)$, proving
\eqref{eq:intensity} and the mean-spacing assertion by telescoping.

Scale each ordinate to $y=\Lambda t$ and put
\[
 w_p=\frac{\lambda_p}{\Lambda}=\frac{\log p}{\sum_{r\in P}\log r},
 \qquad b=\frac{1}{\max_{p\in P}w_p},\qquad
 A_P(x)=\prod_{p\in P}(1-w_p x).
\]
The normalized lattices have respective frequencies $w_p$, summing to one.
The following specialization of the Barra--Gaspard spacing law, as stated by
Marklof and Str\"ombergsson \cite[Theorems 4 and 5]{marklof}, also records the
strict-tail convention at its endpoint.

\begin{proposition}[Gap law for the omitted factors]\label{prop:forward}
For every $x\ge0$ the limit in \eqref{eq:surv} exists, is independent of the
phases, and equals
\begin{equation}\label{eq:law}
 S_P(x)=
 \begin{cases}
  -A_P'(x),&0\le x<b,\\
  0,&x\ge b.
 \end{cases}
\end{equation}
If $p_*$ is the largest prime in $P$, the limiting gap law has an atom at $b$ of mass
\begin{equation}\label{eq:atom}
 w_{p_*}\prod_{r\in P\setminus\{p_*\}}(1-w_r/w_{p_*}).
\end{equation}
\end{proposition}

\begin{proof}
Start at the $n$th point $y=(n+\alpha_p)/w_p$ of the $p$-lattice.
Its next point is $1/w_p$ away. In the $r$-lattice, $r\ne p$, the position
of $y$ within one period is the fractional part of
\[
 n\frac{w_r}{w_p}+\alpha_p\frac{w_r}{w_p}-\alpha_r .
\]
As $n$ increases these positions equidistribute jointly in the unit
$(m-1)$-cube. Indeed, $1$ and all $w_r/w_p$ for $r\ne p$ are rationally
independent. Clearing denominators in a putative relation and multiplying by $\log p$ would
give a nonzero integer relation among logarithms of distinct primes.
Exponentiating rules that out. For each nonzero integer frequency vector the
Fourier average of the sequence is consequently a geometric sum with ratio
different from one, divided by the number of terms, so Weyl's criterion applies.

For $0\le x<b$, the fraction of starting points of type $p$ such that the next
$r$-point lies farther than $x$ tends to $1-w_rx$. Imposing all the $r$-conditions
gives the product of these lengths by joint equidistribution. Their rectangle
has boundary of measure zero. The next $p$-point automatically lies farther
than $x$. Among the first $N$ starting points the proportion from lattice $p$
tends to $w_p$ by the counting formula. Hence
\[
 S_P(x)=\sum_{p\in P}w_p\prod_{r\ne p}(1-w_rx)=-A_P'(x).
\]
For $x\ge b$ every gap is at most $b$, the period of the most frequent lattice.
The left limit at $b$ gives \eqref{eq:atom}. The empty product gives mass one
when $m=1$.
\end{proof}

\section{Arithmetic inversion}
\begin{proof}[Proof of Theorem~\ref{thm:main}]
Each survival function is a polynomial near zero. Since $A_P(0)=1$, its germ
recovers the polynomial
\begin{equation}\label{eq:invert}
 A_P(x)=1-\int_0^x S_P(u)\,du \qquad(0<x<b).
\end{equation}
A polynomial agreeing with these values is unique. Its degree yields $m$, and
its roots are exactly $1/w_p$. This holds even though only the smallest root is
an endpoint of the observed support. If two laws agree near zero, their
sorted normalized weights therefore agree term by term. Write
$p_1<\cdots<p_m$ and $q_1<\cdots<q_m$ for the ordered primes. With
$c=(\sum_i\log q_i)/(\sum_i\log p_i)>0$, we obtain
\[
  q_i=p_i^c\qquad(1\le i\le m).
\]
For $m\ge3$ consider the three rationally independent real numbers
$\log p_1,\log p_2,\log p_3$. If $c$ were irrational, $1,c$ would be
rationally independent as well. The six exponentials theorem then implies
that at least one among
\[
 \exp(\log p_i),\quad \exp(c\log p_i)\qquad (1\le i\le3)
\]
is transcendental \cite{wald}. All six are primes. Hence $c=a/d$ for coprime
positive integers $a,d$. From $q_i^d=p_i^a$, unique factorization gives
$q_i=p_i$ and then $d=a$. Thus $c=1$ and $P=Q$.
\end{proof}

The conclusion determines precisely the set in \eqref{eq:euler}. The forward
law retains no phase information. Powers of primes in $q$ also do not change
that set. For a squarefree principal modulus $q$ with at least three prime
divisors, $P$ is the full prime divisor set, so the short-gap law recovers $q$.

\paragraph{One or two omitted primes.}
For a single prime the normalized law is always the point mass at one.
For two primes $p_1<p_2$, a second pair $q_1<q_2$ gives the same law exactly
when $\log q_1/\log p_1=\log q_2/\log p_2$. A distinct pair would require
an irrational common ratio $c$. The four exponentials conjecture would rule
this out, applied to $1,c$ and $\log p_1,\log p_2$ \cite{wald}.
This gives a conditional extension to two primes. The equivalence just
proved concerns only prime pairs and is a special case of the transcendence
problem.

\section{Approximation by unrelated prime sets}
The inverse theorem uses an exact distribution. Weak convergence of the gap laws means convergence of their survival functions at all continuity points of the limit.
\begin{proposition}\label{prop:instability}
For every finite $P$ with $m\ge2$ there are pairwise disjoint sets $Q_n$ of
$m$ primes, all disjoint from $P$, with $\min Q_n\to\infty$, whose normalized
gap laws converge weakly to the law for $P$. The coefficients of their
short-gap polynomials converge to those for $P$ as well.
\end{proposition}
\begin{proof}
Order the weights $0<w_1<\cdots<w_m$ of $P$. For each sufficiently large $T$
choose a prime $q_i$ between $\lfloor e^{Tw_i}\rfloor$ and
$2\lfloor e^{Tw_i}\rfloor$ by Bertrand's postulate. These intervals are
eventually disjoint as $w_i<w_{i+1}$, and $\log q_i=Tw_i+O(1)$. Thus their
normalized weights are $w_i+O_m(1/T)$ because $\sum_i w_i=1$. Let $T$ tend
to infinity along a sequence growing fast enough that all new primes exceed
the preceding maximum and all primes of $P$. Coefficient convergence follows
from the polynomial product. The cutoffs converge to $b$ and
\eqref{eq:law} gives convergence of the survival functions at every $x\ne b$.
\end{proof}

\section{An example with two different characters}
Take the principal character modulo $30$, and the character modulo $210$ induced
from the quadratic character modulo $7$. Both omit $P=\{2,3,5\}$. Their phases
at these primes are $(0,0,0)$ and $(0,1/2,1/2)$, respectively. Writing
$e_j$ for the elementary symmetric polynomial of degree $j$ in
$\log2/\log30,\log3/\log30,\log5/\log30$, we obtain
\[
 S_P(x)=1-2e_2x+3e_3x^2
 \simeq 1-0.63021791 x+0.09344797x^2,\qquad 0\le x<b.
\]
The cutoff is $b=\log30/\log5\simeq2.11328275$. Approximately
$8.55\%$ of all gaps attain that maximum. In particular, extrapolating the
quadratic beyond $b$ would give incorrect probabilities.

\begin{center}
\small
\begin{tabular}{@{}lrrr@{}}
\toprule
Event & Limiting frequency & Modulo $30$ & Modulo $210$\\
\midrule
$g>0.25$ & 0.84828602 & 0.84828953 & 0.84828640\\
$g>0.75$ & 0.57990106 & 0.57990372 & 0.57990342\\
$g>1.25$ & 0.35824008 & 0.35824309 & 0.35823807\\
$g>1.75$ & 0.18330308 & 0.18330581 & 0.18330698\\
$g>2.00$ & 0.11335608 & 0.11336119 & 0.11335338\\
$g=b$    & 0.08550665 & 0.08551044 & 0.08549535\\
\bottomrule
\end{tabular}
\end{center}
The two samples contain $541{,}314$ and $541{,}316$ complete gaps below height $10^6$. The program \texttt{verify\_gaps.py} computes each ordinate directly and identifies a maximal gap by consecutive points of the $5$-lattice.

\paragraph{Prior work and further questions.}
Barra and Gaspard studied the spacing law of these superposed lattices in
quantum-graph spectra. Marklof and Str\"ombergsson give the exact forward formula
and endpoint atom \cite{marklof}. The product also appears in
Lu et al.\ \cite[eqs.~(6)--(9)]{lu}, following the superposition method of
Berry and Robnik \cite{berry}. The arithmetic contribution here is the
scale-free recovery of omitted primes and its instability under approximation.
A targeted literature search did not establish priority for this combination. For primes bounded by $B$, distinct sets of at least three
primes have distinct gap polynomials and there are only finitely many candidates.
How does the minimum separation of their coefficient vectors vary with $B$?
What finite observation height and precision suffice to distinguish them?

\begin{thebibliography}{9}\footnotesize\setlength{\itemsep}{1.5pt}\setlength{\parskip}{0pt}\raggedright
\bibitem{lean} L. Sernau et al., \emph{GRH for Dirichlet $L$-functions, in Lean 4},
commit \texttt{72fead102f66} (2026).
\href{https://github.com/lsernau/grh-lean/blob/72fead102f663d5c7a407964a13dc0cf9d977207/grh/src/New/Reduction/OneSidedBound.lean}{\texttt{OneSidedBound.lean}}, final change-of-level step.
\bibitem{dlmf} NIST Digital Library of Mathematical Functions,
\href{https://dlmf.nist.gov/25.15}{\S25.15}, equations (25.15.4) and (25.15.5).
\bibitem{ked} K. S. Kedlaya, \emph{Notes on analytic number theory},
\href{https://kskedlaya.org/ant/chap-lfunc.html}{chapter on Dirichlet characters and $L$-functions}.
\bibitem{marklof} J. Marklof and A. Str\"ombergsson, Gaps between logs,
\emph{Bull. Lond. Math. Soc.} \textbf{45} (2013), 1267--1280.
\href{https://doi.org/10.1112/blms/bdt060}{doi:10.1112/blms/bdt060}. See also
F. Barra and P. Gaspard, On the level spacing distribution in quantum graphs,
\emph{J. Statist. Phys.} \textbf{101} (2000), 283--319.
\bibitem{lu} J. Lu, T. Hofmann, U. Kuhl and H.-J. St\"ockmann, Implications of
spectral interlacing for quantum graphs, \emph{Entropy} \textbf{25} (2023), 109.
\href{https://doi.org/10.3390/e25010109}{doi:10.3390/e25010109}.
\bibitem{berry} M. V. Berry and M. Robnik, Semiclassical level spacings when
regular and chaotic orbits coexist, \emph{J. Phys. A} \textbf{17} (1984), 2413--2421.
\href{https://doi.org/10.1088/0305-4470/17/12/013}{doi:10.1088/0305-4470/17/12/013}.
\bibitem{wald} M. Waldschmidt, Further variations on the six exponentials theorem,
\emph{Hardy--Ramanujan J.} \textbf{28} (2005).
\href{https://doi.org/10.46298/hrj.2005.86}{doi:10.46298/hrj.2005.86}.
See also the author's discussion of prime powers in
\href{https://webusers.imj-prg.fr/~michel.waldschmidt/articles/pdf/LIL.pdf}{\emph{Linear Independence of Logarithms of Algebraic Numbers}}, introduction.
\end{thebibliography}
\end{document}
